% !PRINTED 164
\weberpage{164}
%<<CONT
identical vanishing does not take place, but rather a final equation of order $mnp$ appears, then this must also be the
degree of the final equation in the general case.

With this the \gesperrt{Bezout} theorem is proved for the case of
three homogeneous quaternary equations.

That we have restricted ourselves here to three equations
was done merely in the interest of simplicity. The extension required for the case of $n$ equations in $n + 1$ homogeneous unknowns follows of itself, and
may be left to the reader.

With regard to our general problem we wish only to remark further that from these considerations it follows that the
problem of algebra, in so far as it relates to the solution of any
algebraic equations whatever, is thereby reduced
% German Sperrung starts at "ist" and excludes "damit zurückgeführt"; English keeps
% "is thereby reduced" outside since the verb cluster straddles the boundary
\gesperrt{to the treatment of a chain of equations each in
a single unknown.}

\begin{center}
\S.\,51.

\gesperrt{Decomposable and indecomposable functions.}
\end{center}

Among the integral rational functions of several variables, \gesperrt{decomposable and indecomposable} ones must be
distinguished.

An integral rational function of any variables whatever is called \gesperrt{decomposable} if it can be represented as a product of at least
two integral rational functions of the same variables; if it cannot be represented as such a product,
it is called \gesperrt{indecomposable}.

If a function is decomposable, then the degree of each
of the factors is lower than the degree of the function itself. A
\gesperrt{linear} function is therefore always indecomposable, and \gesperrt{every integral
rational function can be decomposed into a finite number
of indecomposable functions.}

We call an integral function $W$ divisible by another $w$
if a third integral function $w'$ exists such that

\[
W = w\,w'
\]

holds. From this it then follows that, if $U$ is an integral function divisible by $w$
and $V$ is an arbitrary integral function, the product $UV$
%CONT>>
